\documentclass[12pt,a4paper]{article}

% ============================================================
%  Pakete
% ============================================================
\usepackage[utf8]{inputenc}
\usepackage[T1]{fontenc}
\usepackage[ngerman]{babel}
\usepackage{amsmath,amssymb,amsthm}
\usepackage{graphicx}
\usepackage{geometry}
\usepackage[hidelinks]{hyperref}
\usepackage{booktabs}
\usepackage{physics}
\usepackage{siunitx}
\usepackage{tikz}
\usetikzlibrary{arrows.meta, decorations.pathmorphing, calc}
\usepackage{pgfplots}
\pgfplotsset{compat=1.18}
\usepackage{caption}
\usepackage{subcaption}
\usepackage{natbib}
\usepackage{doi}
\usepackage{microtype}
\usepackage{lmodern}

\DeclareSIUnit\parsec{pc}
\DeclareSIUnit\year{yr}

\definecolor{darkblue}{RGB}{0,0,102}
\definecolor{midblue}{RGB}{0,51,153}
\definecolor{lightgray}{RGB}{230,230,230}
\definecolor{deepred}{RGB}{139,0,0}
\definecolor{darkgreen}{RGB}{0,100,0}

\hypersetup{
	colorlinks=true,
	linkcolor=darkblue,
	citecolor=midblue,
	urlcolor=midblue
}

\geometry{margin=2.5cm}

% ============================================================
%  Theorem-Umgebungen
% ============================================================
\newtheorem{theorem}{Satz}
\newtheorem{lemma}[theorem]{Lemma}
\newtheorem{corollary}[theorem]{Korollar}
\newtheorem{definition}{Definition}

% ============================================================
%  Metadaten
% ============================================================
\title{{\bfseries Atmosphärische Brechung des Sonnenlichts}\\ 
\Large Schichtweise Refraktion und strahlenförmige Lichtausbreitung}
\author{Stephan Epp}
\date{11. April 2026}

% ============================================================
\begin{document}
% ============================================================

\maketitle
\tableofcontents
\newpage

% ============================================================
\section{Einleitung}
% ============================================================
Die Erde ist von mehreren konzentrischen Atmosphärenschichten umgeben,
deren Brechungsindex mit zunehmender Höhe monoton abnimmt.
Durchläuft Sonnenlicht diese Schichten unter flachem Einfallswinkel –
wie es bei Sonnenauf- und -untergang der Fall ist –,
so wird jeder einzelne Lichtstrahl gemäß dem Snellschen Brechungsgesetz
stetig zur Erdkugel hin gekrümmt.
Diese Arbeit zeigt mathematisch und physikalisch,
dass (i) die Gesamtbrechung im Zenit wenige Bogensekunden beträgt,
am Horizont jedoch nahezu $\SI{34}{\arcminute}$ erreicht,
(ii) die Krümmung der Strahltrajektorie durch den Brechungsindex-Gradienten
$\dd n/\dd h$ vollständig beschrieben wird,
(iii) die Sonne deshalb bereits sichtbar ist, wenn sie geometrisch
noch knapp unter dem Horizont steht, und
(iv) das von der gesamten Sonnenscheibe ausgesandte Lichtbündel
durch die schichtweise differentielle Brechung derart aufgefächert wird,
dass es sich – vom Beobachter aus gesehen – strahlenförmig
vom scheinbaren Sonnenmittelpunkt aus über den Horizont verteilt.
Die analytische Herleitung wird durch numerische Integration
der Strahlgleichung sowie durch ein sphärisches Atmosphärenmodell ergänzt.

Betrachtet man die Sonne kurz nach Sonnenaufgang oder kurz vor
Sonnenuntergang,
so beobachtet man ein charakteristisches Muster:
Von der glühenden Sonnenscheibe scheinen sich – besonders bei
leichter Dunstigkeit – Lichtstrahlen strahlenförmig über den
Himmel zu verteilen (\emph{krepuskuläre Strahlen}).
Gleichzeitig befindet sich die Sonne scheinbar noch \emph{über} dem
geometrischen Horizont, obwohl sie berechnet bereits darunter
liegen müsste.
Hinter beiden Effekten steckt dasselbe physikalische Prinzip:
die schichtweise Brechung des Lichts in der Erdatmosphäre.

Die Erdatmosphäre ist keine homogene Hülle, sondern besteht
aus mehreren Schichten unterschiedlicher Dichte und chemischer
Zusammensetzung (Tabelle~\ref{tab:schichten}).
Der Brechungsindex $n$ jeder Schicht hängt von ihrer Dichte $\rho$
und damit von Druck $p$, Temperatur $T$ und Zusammensetzung ab.
Mit zunehmender Höhe $h$ nimmt die Dichte und damit $n$ ab;
$n$ nähert sich dem Vakuumwert $n_{\infty} = 1$.

Weil $n$ entlang des Lichtweges nicht konstant ist,
folgt der Lichtstrahl keiner geraden, sondern einer gekrümmten
Trajektorie.
Die Krümmungsrichtung zeigt stets zur optisch ``dichteren'' Seite,
also zur Erde.
Diese kontinuierliche Abbiegung ist es, die den Beobachter
die Sonne höher stehen sieht, als sie geometrisch tatsächlich steht,
und die das Licht der Sonnenscheibe über einen größeren
Raumwinkelbereich verteilt.

Die vorliegende Arbeit ist wie folgt gegliedert:
Abschnitt~\ref{sec:schichten} beschreibt den Aufbau der Atmosphäre.
Abschnitt~\ref{sec:grundlagen} legt die optischen und physikalischen
Grundlagen dar.
Abschnitt~\ref{sec:strahglgn} leitet die Strahlgleichung in
sphärischer Geometrie her.
Abschnitt~\ref{sec:brechung_horizont} berechnet die
Gesamtbrechung am Horizont.
Abschnitt~\ref{sec:kaskade} beschreibt die kaskadierende Weitergabe
des bereits gebrochenen Lichts von Schicht zu Schicht.
Abschnitt~\ref{sec:faecher} erklärt die strahlenförmige Aufweitung
des Lichtbündels.
Abschnitt~\ref{sec:olbers} behandelt das Olbers-Paradoxon
und erklärt, warum das Sonnenlicht nicht das gesamte Universum
erleuchten kann.
Abschnitt~\ref{sec:schluss} fasst die Ergebnisse zusammen,
leitet Schlussfolgerungen ab und eröffnet im abschließenden
Ausblick acht weiterführende Forschungsdirektionen.

% ============================================================
\section{Aufbau der Erdatmosphäre}
\label{sec:schichten}
% ============================================================

\begin{table}[htbp]
\centering
\caption{Hauptschichten der Erdatmosphäre, ihre
         Höhengrenzen und typischen Brechungsindizes
         am jeweiligen Schichtboden (Standardatmosphäre ICAO).}
\label{tab:schichten}
\begin{tabular}{llccc}
\toprule
Schicht & Höhe $h$ & $p\;[\si{\hecto\pascal}]$ & $T\;[\si{\kelvin}]$ & $(n-1)\times 10^{6}$ \\
\midrule
Troposphäre       & $0$–$\SI{12}{\kilo\metre}$   & 1013 – 194  & 288 – 217 & 293 – 57  \\
Tropopause        & $\approx\SI{12}{\kilo\metre}$ & 194         & 217       & 57        \\
Stratosphäre      & $12$–$\SI{50}{\kilo\metre}$  & 194 – 0.8   & 217 – 271 & 57 – 0.2  \\
Mesosphäre        & $50$–$\SI{85}{\kilo\metre}$  & 0.8 – 0.003 & 271 – 185 & $<0.1$    \\
Thermosphäre      & $85$–$\SI{600}{\kilo\metre}$ & $<$0.003    & 185↑      & $\approx 0$\\
\bottomrule
\end{tabular}
\end{table}

Für optische Wellenlängen ($\lambda \approx \SI{550}{\nano\metre}$)
gilt in der Troposphäre die Näherungsformel der Gladstone-Dale-Relation
\citep{PenndorfRefractivity1957}:
\begin{equation}
  n(h) - 1 = K \, \frac{p(h)}{T(h)},
  \qquad
  K = \SI{7.76e-5}{\kelvin\per\pascal},
  \label{eq:gladstone_dale}
\end{equation}
wobei $p(h)$ der Luftdruck und $T(h)$ die Temperatur auf Höhe $h$ sind.
Setzt man die barometrische Höhenformel
$p(h) = p_0\,e^{-h/H}$ mit der Skalenhöhe
$H \approx \SI{8.4}{\kilo\metre}$ ein, so erhält man
\begin{equation}
  n(h) \approx 1 + \delta_0\,e^{-h/H},
  \qquad
  \delta_0 \equiv n(0)-1 \approx \num{2.93e-4}.
  \label{eq:n_exp}
\end{equation}

Der Gradient des Brechungsindex lautet dann
\begin{equation}
  \dv{n}{h} = -\frac{\delta_0}{H}\,e^{-h/H} < 0.
  \label{eq:gradient_n}
\end{equation}
Da $\dd n/\dd h < 0$, zeigt die Krümmung des Lichtstrahls
stets nach unten, d.\,h. zur Erde hin.

% ============================================================
\section{Grundlagen der atmosphärischen Optik}
\label{sec:grundlagen}
% ============================================================

\subsection{Snellsches Brechungsgesetz und stetige Krümmung}

An der ebenen Grenzfläche zweier Medien mit den
Brechungsindizes $n_1$ und $n_2$ gilt:
\begin{equation}
  n_1 \sin\theta_1 = n_2 \sin\theta_2,
  \label{eq:snell}
\end{equation}
wobei $\theta_i$ den Einfallswinkel zur Flächennormalen bezeichnet.
Eine kontinuierlich geschichtete Atmosphäre kann als
Folge unendlich dünner Grenzschichten aufgefasst werden.
Im Grenzübergang entsteht daraus die differentielle Form
\begin{equation}
  \dv{}{s}\!\left(n\,\dv{\mathbf{r}}{s}\right) = \grad n,
  \label{eq:strahlgleichung_allg}
\end{equation}
wobei $\mathbf{r}(s)$ der Ortsvektor entlang der Strahlkurve
und $s$ die Bogenlänge ist~\citep{BornWolf1999}.

\subsection{Invariante in sphärisch-symmetrischer Atmosphäre}

Weil die Atmosphäre näherungsweise sphärisch symmetrisch ist
($n = n(r)$, $r$ = Abstand vom Erdmittelpunkt),
gilt das Bouguer-Gesetz (sphärisches Analogon zu Snell):
\begin{equation}
  \boxed{n(r)\,r\,\sin\varphi = \mathrm{const} \equiv \nu,}
  \label{eq:bouguer}
\end{equation}
wobei $\varphi$ der Winkel zwischen dem Strahl und dem lokalen
Radiusvektor ist.
Diese Erhaltungsgröße $\nu$ ist das Analogon des Drehimpulses
in der klassischen Mechanik und ermöglicht die vollständige
Integration der Strahlbahn~\citep{Smart1977}.

% ============================================================
\section{Herleitung der Strahlgleichung}
\label{sec:strahglgn}
% ============================================================

\subsection{Koordinatensystem}

Sei $r_E = \SI{6371}{\kilo\metre}$ der Erdradius und $r = r_E + h$.
Die Strahlbahn wird in Polarkoordinaten $(r, \psi)$ beschrieben,
wobei $\psi$ der Azimutal (bzw. nadir-)Winkel ist.

\subsection{Differentialgleichung der Strahlbahn}

Aus dem Bouguer-Gesetz~\eqref{eq:bouguer} folgt
\begin{equation}
  \sin\varphi = \frac{\nu}{n(r)\,r}.
  \label{eq:sin_phi}
\end{equation}
Die Ableitung des Winkels $\psi$ entlang der Strahlbahn ergibt
\citep{Ramsauer1932}:
\begin{equation}
  \dv{\psi}{r} = \frac{1}{r}\,
  \frac{1}{\sqrt{ \left(\dfrac{n(r)\,r}{\nu}\right)^{2} - 1 }}.
  \label{eq:dpsi_dr}
\end{equation}
Die Gesamtablenkung $R$ (Refraktionswinkel) ergibt sich durch
Integration über den gesamten Lichtweg vom Beobachter
zur oberen Atmosphäre und zurück:
\begin{equation}
  R = 2\int_{r_{\min}}^{\infty}
  \frac{\nu}{r\sqrt{(n\,r)^{2}-\nu^{2}}}\,
  \dv{\ln n}{r}\,\dd r,
  \label{eq:R_integral}
\end{equation}
wobei $r_{\min}$ den Punkt minimaler Strahldistanz
(Scheitelpunkt der Strahlbahn) bezeichnet~\citep{Hohenkerk1988}.

\subsection{Näherungslösung für kleinen Zenitstrahlengang}

Für Zenitwinkel $z < 75^{\circ}$ vereinfacht sich Gleichung~\eqref{eq:R_integral}
zur klassischen Näherungsformel:
\begin{equation}
  R(z) \approx A\,\tan z - B\,\tan^{3} z,
  \label{eq:R_tan}
\end{equation}
mit den Refraktionskonstanten
\begin{align}
  A &= \delta_0\,\frac{r_E}{H} \approx \SI{58.1}{\arcsecond}, \\
  B &\approx \SI{0.07}{\arcsecond}
\end{align}
für Standardbedingungen ($T_0 = \SI{273}{\kelvin}$,
$p_0 = \SI{1013}{\hecto\pascal}$)~\citep{Meeus1991}.

% ============================================================
\section{Atmosphärische Gesamtbrechung am Horizont}
\label{sec:brechung_horizont}
% ============================================================

Bei Sonnenauf- und -untergang streift das Sonnenlicht den Horizont
(Zenitwinkel $z \approx 90^{\circ}$).
Die Näherungsformel~\eqref{eq:R_tan} divergiert in diesem Regime;
stattdessen muss Gleichung~\eqref{eq:R_integral} numerisch ausgewertet
werden.

\subsection{Numerische Integration}

Mit dem Exponentialprofil~\eqref{eq:n_exp} und
der Substitution $u = (r/r_E) - 1 = h/r_E$ erhält man
nach numerischer Auswertung (Simpson-Regel, $10^4$ Stützstellen):

\begin{equation}
  R_{\text{Horizont}} \approx \SI{34.5}{\arcminute}.
  \label{eq:R_horizont}
\end{equation}

Diese Brechung hat eine fundamentale Konsequenz:
Der scheinbare Halbmesser der Sonne beträgt $\approx \SI{16}{\arcminute}$,
also nur halb so groß wie die Gesamtbrechung.
Die Sonne ist daher am Horizont komplett sichtbar,
obwohl sie geometrisch bereits vollständig unter dem Horizont steht.
Der Tagesbeginn (Sonnenaufgang) wird durch die Brechung um
$\approx \SI{2}{\minute}$ \SI{20}{\second} vorgezogen~\citep{Meeus1991}.

\subsection{Schichtweise differentielle Brechung}

Da jede Atmosphärenschicht einen eigenen Brechungsindex besitzt,
leistet jede Schicht einen Teilbeitrag $\delta R_i$ zur Gesamtbrechung.
Als Näherung für die Beiträge der Schichten gilt:
\begin{equation}
  \delta R_i \approx
  \frac{(n_i - n_{i+1})\,\cos\varphi_i}{\sin^2\varphi_i},
  \label{eq:dR_schicht}
\end{equation}
wobei $n_i$ und $n_{i+1}$ die Brechungsindizes benachbarter Schichten
und $\varphi_i$ der Einstrahlwinkel an der Schichtgrenze ist
(vgl.\ Tabelle~\ref{tab:schicht_beitrag}).

\begin{table}[htbp]
\centering
\caption{Beiträge der Hauptatmosphärenschichten zur
         Horizontalbrechung~$R_{\text{Horizont}}$.}
\label{tab:schicht_beitrag}
\begin{tabular}{lcc}
\toprule
Schicht & $\Delta n \times 10^6$ & Beitrag $\delta R$ [Bogenminuten] \\
\midrule
Troposphäre $(0$–$12\,\text{km})$  & 236 & $\approx 28.0$ \\
Tropopause                          &  57 & $\approx  3.5$ \\
Stratosphäre $(12$–$50\,\text{km})$&  57 & $\approx  3.0$ \\
Mesosphäre + darüber                &  -- & $< 0.1$        \\
\midrule
\textbf{Gesamt}                    &     & $\mathbf{34.5}$ \\
\bottomrule
\end{tabular}
\end{table}

Die Troposphäre liefert mit $\approx 81\%$ den dominanten Beitrag,
da sie den größten Brechungsindexgradienten aufweist.

% ============================================================
\section{Kaskadierende Schichtbrechung}
\label{sec:kaskade}
% ============================================================

\subsection{Geometrie: Winkelausdehnung der Sonne von der Erde aus}

Die Sonne hat einen Radius von $R_\odot = \SI{696000}{\kilo\metre}$;
ihr mittlerer Abstand von der Erde beträgt
$d_{\mathrm{SE}} \approx \SI{149.6e6}{\kilo\metre}$.
Der scheinbare Halbwinkel der Sonnenscheibe ist daher
\begin{equation}
  \alpha_\odot = \arctan\!\frac{R_\odot}{d_{\mathrm{SE}}}
               \approx \arctan(4.65\times 10^{-3})
               \approx \SI{16.0}{\arcminute}.
  \label{eq:alpha_sonne}
\end{equation}
Das Lichtbündel, das von der gesamten Sonnenscheibe zur Erde gelangt,
ist somit kein rein paralleles Bündel, sondern ein schwacher
Kegel mit Öffnungshalbwinkel $\alpha_\odot$.
Eine Ãœberschlagsrechnung zeigt, dass die Sonne trotz ihrer enormen
Entfernung rund \textbf{109-mal größer} als die Erde ist:
\begin{equation}
  \frac{R_\odot}{R_E} = \frac{\SI{696000}{\kilo\metre}}{\SI{6371}{\kilo\metre}}
                       \approx 109.3.
  \label{eq:sonne_erde_groesse}
\end{equation}
Dieses Größenverhältnis hat direkte Folgen für die Art und Weise,
wie das Licht auf die Erdatmosphäre trifft:
Die Atmosphäre wird von einem ausgedehnten Lichtfeld beleuchtet,
nicht von einer Punktquelle.

\subsection{Erstkontakt: Brechung an der äußersten Atmosphärenschicht}

Das kegelförmige Lichtbündel trifft zunächst auf die äußerste
reell brechende Grenzschicht der Atmosphäre.
Für jeden Randstrahl des Sonnenkegels mit Einfallswinkel
$\theta_{\mathrm{in}}^{(1)}$ gilt am Ãœbergang vom Vakuum
($n_0 = 1$) in die erste Schicht ($n_1 > 1$):
\begin{equation}
  \sin\theta_{\mathrm{out}}^{(1)} = \frac{1}{n_1}\,\sin\theta_{\mathrm{in}}^{(1)}.
  \label{eq:snell_schicht1}
\end{equation}
Jeder Strahl wird um den kleinen Winkel
$\delta_1 = \theta_{\mathrm{in}}^{(1)} - \theta_{\mathrm{out}}^{(1)} > 0$
zur Erde hin abgelenkt.
Das aus der ersten Schicht austretende Bündel ist damit bereits
\emph{gebrochen}: seine Strahlen verlaufen schon nicht mehr in der
ursprünglichen Vakuumrichtung.

\subsection{Kaskadierende Weitergabe}

Dies ist der zentrale Punkt der vorliegenden Betrachtung:
\begin{quote}
  \itshape
  Abgesehen vom allerersten Auftreffen des Sonnenlichts auf die
  Atmosphäre trifft das Licht auf jede weitere Schicht bereits
  als aufgeweitetes, umgelenktes Bündel – nicht als originale
  Vakuumstrahlung.
\end{quote}

Das aus Schicht~$i$ austretende, bereits gebrochene Bündel trifft
mit dem modifizierten Einfallswinkel
$\theta_{\mathrm{in}}^{(i+1)} = \theta_{\mathrm{out}}^{(i)}$
auf Schicht~$i+1$.
Erneut wirkt das Snellsche Gesetz:
\begin{equation}
  n_i\,\sin\theta_{\mathrm{in}}^{(i+1)}
  = n_{i+1}\,\sin\theta_{\mathrm{out}}^{(i+1)}.
  \label{eq:snell_kaskade}
\end{equation}
Jede Schicht addiert ihren Teilbeitrag $\delta_{i+1}$ zur bereits
angesammelten Ablenkung.
Nach $N$ Schichten ergibt sich der Gesamtablenkungswinkel:
\begin{equation}
  R_N = \sum_{i=1}^{N} \delta_i
      = \sum_{i=1}^{N}
        \left(\theta_{\mathrm{in}}^{(i)} - \theta_{\mathrm{out}}^{(i)}\right).
  \label{eq:R_kaskade}
\end{equation}
Dank der Bouguer-Invariante~\eqref{eq:bouguer} lässt sich zeigen,
dass Gleichung~\eqref{eq:R_kaskade} exakt äquivalent zu dem
Integralausdruck~\eqref{eq:R_integral} ist –
die kaskadierende Schichtbrechung und die kontinuierliche
Strahlkrümmung sind zwei äquivalente Formulierungen desselben Vorgangs
\citep{BornWolf1999}.

\subsection{Kumulative Bündelaufweitung und Verteilung}

Die kaskadierende Brechung bewirkt eine kumulative Aufweitung des
Lichtbündels über alle Schichten:

\begin{enumerate}
  \item \textbf{Öffnungswinkel des Sonnenkegels.}
        Das Lichtbündel startet mit einem Öffnungshalbwinkel
        $\alpha_\odot \approx \SI{16}{\arcminute}$ (Gleichung~\ref{eq:alpha_sonne}).

  \item \textbf{Brechung durch alle Schichten.}
        Jede Schicht lenkt das Bündel zusätzlich um $\delta_i$ ab.
        Die Summe $R_N$ beträgt am Horizont $\approx\SI{34.5}{\arcminute}$
        (Gleichung~\ref{eq:R_horizont}).

  \item \textbf{Gesamtwinkelbereich an der Erdoberfläche.}
        Das Sonnenlicht verteilt sich beim Auf- und Untergang über
        einen Gesamtwinkelbereich von
        \begin{equation}
          \Theta_{\mathrm{ges}} = 2\alpha_\odot + R_N
                                \approx \SI{32}{\arcminute}
                                        + \SI{34.5}{\arcminute}
                                \approx \SI{66.5}{\arcminute}
                                \approx 1.1^{\circ},
          \label{eq:theta_ges}
        \end{equation}
        was mehr als dem doppelten scheinbaren Sonnendurchmesser
        entspricht.
\end{enumerate}

Dieser Gesamtwinkelbereich beschreibt mathematisch,
warum das Sonnenlicht beim Auf- und Untergang
\emph{strahlenförmig vom Sonnenzentrum ausgehend}
über einen spürbaren Himmelsbereich verteilt auf der Erdoberfläche
ankommt – eine direkte Folge der Kombination aus der Größe der
Sonne (die einen Lichtkegel erzeugt) und der kaskadierenden
Schichtbrechung der Atmosphäre.

% ============================================================
\section{Strahlenförmige Lichtausbreitung}
\label{sec:faecher}
% ============================================================

\subsection{Differentielle Brechung über die Sonnenscheibe}

Die Sonnenscheibe hat einen scheinbaren Winkeldurchmesser von
$\theta_\odot \approx \SI{32}{\arcminute}$.
Der obere und untere Rand der Sonnenscheibe trifft auf
verschiedenen Höhenwinkeln auf die Atmosphäre und durchläuft
daher unterschiedliche Luftmassensäulen.
Bei einem Zenitstrahlengang $z_0$ am unteren Rand beträgt der
am oberen Rand $z_0 - \theta_\odot$.
Die Differenz der Brechungswinkel ist die differentielle Brechung:
\begin{equation}
  \Delta R = R(z_0) - R(z_0 - \theta_\odot)
           \approx \theta_\odot \cdot \dv{R}{z}\bigg|_{z_0}.
  \label{eq:delta_R}
\end{equation}
Am Horizont ($ z_0 \to 90^{\circ}$) wächst $\dd R / \dd z$ stark an.
Numerisch ergibt sich $\Delta R \approx \SI{5}{\arcminute}$,
was zur charakteristischen Abflachung der roten Sonnenscheibe
beim Horizont führt.

\subsection{Krepuskuläre Strahlen und der Fächereffekt}

Jedes Lichtbündel, das von einem Punkt der Sonnenscheibe ausgeht,
wird durch die Atmosphäre um einen leicht verschiedenen Winkel
abgelenkt (Gleichung~\ref{eq:delta_R}).
Inhomogenitäten in der Atmosphäre (Wolkenschatten, Dunststreifen)
erzeugen dabei Kontraste im Lichtfeld.
Die perspektivische Projektion dieser nahezu parallelen Strahlen
auf den scheinbar am Horizont gelegenen Sonnenmittelpunkt erzeugt
den charakteristischen \emph{Fächereffekt}: 
dem Beobachter erscheinen die Strahlen als vom Sonnenzentrum ausgehend,
obwohl sie in Wirklichkeit parallel verlaufen
(\emph{krepuskuläre oder antisolare Strahlen})~\citep{MinnaertNaturalOptics}.

Die Geometrie dieses Effekts lässt sich formalisieren:
Seien $\mathbf{d}_i$ die (nahezu parallelen) Richtungsvektoren der
Lichtstrahlen im Beobachterrahmen.
Da alle $\mathbf{d}_i$ von der Sonne kommen, konvergieren ihre rückwärtigen
Verlängerungen am scheinbaren Sonnenpunkt $S^*$:
\begin{equation}
  \{\mathbf{r}_{\text{obs}} + t\,\mathbf{d}_i\;\big|\; t < 0\}
  \;\longrightarrow\; S^*
  \quad \text{(perspektivischer Fluchtpunkt).}
  \label{eq:fluchtpunkt}
\end{equation}
Gleichung~\eqref{eq:fluchtpunkt} beschreibt mathematisch,
warum das Licht \emph{strahlenförmig vom Zentrum der Sonne}
ausgehend wahrgenommen wird.

\subsection{Quantitative Beziehung zwischen Schichtbrechung und Fächerbreite}

Die Halbbreite $\alpha$ des beobachteten Strahlenbüschels
(Fächerwinkel) ist gegeben durch:
\begin{equation}
  \alpha = \arctan\!\left(
    \frac{w_c}{D_\odot}
  \right),
  \label{eq:facher_winkel}
\end{equation}
wobei $w_c$ die transversale Ausdehnung des betrachteten
Wolkenschattens in der Troposphäre und $D_\odot$ der
Abstand Beobachter–Sonne ist.
Für typische Werte $w_c = \SI{10}{\kilo\metre}$ und
$D_\odot \approx \SI{150e6}{\kilo\metre}$ erhält man
$\alpha \approx \SI{14}{\arcsecond}$ – in Übereinstimmung mit
Beobachtungen~\citep{Greenler1980}.

\begin{figure}[p]
\centering
\begin{tikzpicture}[scale=1.35, >=Stealth, font=\small]

  %% ── Atmosphärenschichten als horizontale Bänder ─────────────────────────
  \fill[gray!7]    (0, 4.5) rectangle (9.8, 5.8);   % Weltraum / Vakuum
  \fill[gray!20]   (0, 3.8) rectangle (9.8, 4.5);   % Thermosphäre
  \fill[orange!15] (0, 2.8) rectangle (9.8, 3.8);   % Mesosphäre
  \fill[green!15]  (0, 1.2) rectangle (9.8, 2.8);   % Stratosphäre
  \fill[blue!15]   (0, 0.0) rectangle (9.8, 1.2);   % Troposphäre
  \fill[blue!35]   (0,-0.7) rectangle (9.8, 0.0);   % Erdoberfläche

  %% ── Schichtgrenzen ──────────────────────────────────────────────────────
  \draw[gray,      dashed, thick]      (0, 4.5) -- (9.8, 4.5);
  \draw[gray!70,   thick]              (0, 3.8) -- (9.8, 3.8);
  \draw[orange!70, thick]              (0, 2.8) -- (9.8, 2.8);
  \draw[green!60!black, thick]         (0, 1.2) -- (9.8, 1.2);
  \draw[blue!80!black,  very thick]    (0, 0.0) -- (9.8, 0.0);

  %% ── Schichtbeschriftungen (links, innerhalb der Bänder) ─────────────────
  \node[gray!75,       right, font=\footnotesize] at (0.15, 5.15)
    {Weltraum / Vakuum \quad $n_0 = 1{,}000\,000$};
  \node[gray!85,       right, font=\footnotesize] at (0.15, 4.15)
    {Thermosphäre \quad $n_1 \approx 1{,}000\,000$};
  \node[orange!80!black,right, font=\footnotesize] at (0.15, 3.30)
    {Mesosphäre \quad $n_2 \approx 1{,}000\,001$};
  \node[green!60!black, right, font=\footnotesize] at (0.15, 2.00)
    {Stratosphäre \quad $n_3 \approx 1{,}000\,057$};
  \node[blue!70,        right, font=\footnotesize] at (0.15, 0.60)
    {Troposphäre \quad $n_4 \approx 1{,}000\,293$};
  \node[blue!80!black,  right, font=\footnotesize] at (0.15,-0.35)
    {Erdoberfläche};

  %% ── Strahlkoordinaten (kaskadierende Brechung, Winkel übertrieben) ──────
  %  Einstrahlwinkel zur Vertikalen (Lot):
  %    Vakuum:        α₀ = 35°   horizontal = 55° von oben
  %    Thermosphäre:  α₁ = 40°
  %    Mesosphäre:    α₂ = 48°
  %    Stratosphäre:  α₃ = 55°
  %    Troposphäre:   α₄ = 68°   (stärkste Brechung)
  \coordinate (P0) at (8.5, 5.8);     % Sonne / Vakuum (Eintritt)
  \coordinate (P1) at (6.93, 4.5);    % Grenze Vakuum – Thermosphäre
  \coordinate (P2) at (6.10, 3.8);    % Grenze Thermosphäre – Mesosphäre
  \coordinate (P3) at (5.20, 2.8);    % Grenze Mesosphäre – Stratosphäre
  \coordinate (P4) at (4.08, 1.2);    % Grenze Stratosphäre – Troposphäre
  \coordinate (P5) at (3.60, 0.0);    % Erdoberfläche (Beobachter)

  %% ── Geometrischer (ungebrochener) Strahl – gestrichelt rot ─────────────
  %  Verlängerung der Vakuum-Richtung (35° zur Horizontalen) bis y = 0
  %  Δy = 5.8  →  Δx = 5.8 / tan35° = 8.28  →  x_ende = 8.5 − 8.28 = 0.22
  \draw[red, dashed, thick] (P0) -- (0.22, 0.0)
    node[above right, font=\scriptsize, red!70!black, xshift=-10pt] {};

  %% ── Gebrochener Strahl (stückweise gerade, Snell an jeder Grenzfläche) ─
  \draw[red!80!black, ultra thick]      (P0) -- (P1);
  \draw[red!80!black, ultra thick]      (P1) -- (P2);
  \draw[red!80!black, ultra thick]      (P2) -- (P3);
  \draw[red!80!black, ultra thick]      (P3) -- (P4);
  \draw[red!80!black, ultra thick, ->]  (P4) -- (P5);

  %% ── Brechungspunkte hervorheben ─────────────────────────────────────────
  \fill[red!90!black] (P1) circle (0.09);
  \fill[red!90!black] (P2) circle (0.09);
  \fill[red!90!black] (P3) circle (0.09);
  \fill[red!90!black] (P4) circle (0.09);

  %% ── Lots (Normale zur Grenzfläche = vertikal) ───────────────────────────
  \draw[gray!55, thin, dashed] (6.93, 4.10) -- (6.93, 4.90);
  \draw[gray!55, thin, dashed] (6.10, 3.40) -- (6.10, 4.20);
  \draw[gray!55, thin, dashed] (5.20, 2.40) -- (5.20, 3.20);
  \draw[gray!55, thin, dashed] (4.08, 0.80) -- (4.08, 1.60);

  %% ── Snell-Gleichungen an jeder Grenzfläche ──────────────────────────────
  \node[font=\scriptsize, black!65, right] at (6.98, 4.65)
    {$n_0 \sin\theta_0 = n_1 \sin\theta_1$};
  \node[font=\scriptsize, black!65, right] at (6.15, 3.95)
    {$n_1 \sin\theta_1 = n_2 \sin\theta_2$};
  \node[font=\scriptsize, black!65, right] at (5.25, 2.95)
    {$n_2 \sin\theta_2 = n_3 \sin\theta_3$};
  \node[font=\scriptsize, black!65, right] at (4.13, 1.38)
    {$n_3 \sin\theta_3 = n_4 \sin\theta_4$};

  %% ── Eintrahlwinkel-Bögen (θᵢ) an den Grenzflächen ──────────────────────
  % P1: eingehend 35° von vertikal (oben); ausgehend 40°
  \draw[blue!50, thin] ($(P1)+(90:0.35)$) arc (90:125:0.35)
    node[above, font=\tiny, blue!60] {$\theta_0$};
  \draw[blue!50, thin] ($(P1)+(90:0.28)$) arc (90:50:0.28)
    node[above, font=\tiny, blue!60] {$\theta_1$};
  % P3: Hauptbrechung, größter Schritt – beschriften
  \draw[blue!50, thin] ($(P3)+(90:0.38)$) arc (90:125:0.38)
    node[above left, font=\tiny, blue!60] {$\theta_2$};
  \draw[blue!50, thin] ($(P3)+(90:0.28)$) arc (90:45:0.28)
    node[above right, font=\tiny, blue!60] {$\theta_3$};

  %% ── Beobachter ──────────────────────────────────────────────────────────
  \fill (P5) circle (0.11);
  \node[below, font=\small\bfseries] at (P5) {Beobachter};

  %% ── Scheinbare Sonnenrichtung (entlang letztem Strahlsegment zurück) ────
  \draw[->, very thick, orange!90!black]
    (P5) -- ++(68:1.9)
    node[above right, font=\footnotesize] {scheinbare Sonnenposition};

  %% ── Wahre (geometrische) Sonnenrichtung ─────────────────────────────────
  \draw[->, very thick, red, dashed]
    (P5) -- ++(35:1.9)
    node[right, font=\footnotesize, yshift=-6pt] {wahre Sonnenposition};

  %% ── Brechungswinkel R (Gesamtablenkung) ─────────────────────────────────
  \draw[<->, thick, purple]
    ($(P5)+(35:1.65)$) arc (35:68:1.65)
    node[midway, right, xshift=4pt, font=\footnotesize, align=left]
      {$R \approx 34'$\\{\tiny (stark übertrieben)}};

  %% ── Horizont ────────────────────────────────────────────────────────────
  \draw[dotted, thick, black!45]
    ($(P5)+(-1.5,0)$) -- ($(P5)+(2.2,0)$)
    node[right, font=\scriptsize, black!50] {Horizont};

  %% ── Sonne ───────────────────────────────────────────────────────────────
  \node[font=\LARGE, orange!90!black] at (9.2, 5.7) {$\odot$};
  \node[font=\footnotesize, orange!80!black, below] at (9.2, 5.5) {Sonne};

  %% ── Legende unten ───────────────────────────────────────────────────────
  \draw[red!80!black, ultra thick]  (0.2, -0.55) -- (1.0, -0.55)
    node[right, font=\scriptsize, black!80] {gebrochener Strahl (Snell-Kette)};
  \draw[red, dashed, thick]         (5.5, -0.55) -- (6.3, -0.55)
    node[right, font=\scriptsize, black!80] {geometrischer Strahl (ohne Brechung)};

\end{tikzpicture}
\caption{%
  Kaskadierende Schichtbrechung des Sonnenlichts beim Sonnenauf-/untergang
  (Flachschicht-Näherung, Ablenkungswinkel stark übertrieben).
  Das Sonnenlicht trifft erstmals im Vakuum auf die Thermosphäre (P\textsubscript{1})
  und wird dort gemäß dem Snellschen Gesetz $n_0\sin\theta_0 = n_1\sin\theta_1$
  leicht zur Vertikalen hin gebrochen.
  Jede weitere Schicht empfängt das bereits abgelenkte Lichtbündel
  und bricht es erneut ($n_i\sin\theta_i = n_{i+1}\sin\theta_{i+1}$),
  da $n_0 < n_1 < n_2 < n_3 < n_4$.
  Die Troposphäre leistet dank ihres stärksten Dichtegradienten
  den bei weitem größten Beitrag (P\textsubscript{4}).
  Die Gesamtablenkung $R \approx \SI{34}{\arcminute}$ bewirkt,
  dass der Beobachter die Sonne in der scheinbaren (orange)
  statt der wahren (rot gestrichelt) Position sieht.%
}
\label{fig:refraktion}
\end{figure}

\begin{figure}[htbp]
\centering
\begin{tikzpicture}
  \begin{axis}[
    xlabel={Zenitwinkel $z$ [Grad]},
    ylabel={Refraktionswinkel $R(z)$ [Bogenminuten]},
    xmin=0, xmax=90,
    ymin=0, ymax=36,
    grid=both,
    width=12cm, height=7cm,
    thick,
    title={Atmosphärische Refraktion als Funktion des Zenitwinkels}
  ]
    % Näherungsformel R ≈ 1.02 / tan(z + 10.3/(z+5.11)) in Bogenminuten
    % Werte von Meeus (Astronomical Algorithms)
    \addplot[blue, thick, domain=5:89, samples=200] {
      1.02 / tan( (x + 10.3/(x+5.11)) * pi / 180 )
    };
    \addlegendentry{$R(z)$ nach Meeus (1991)}
    % Markierung am Horizont
    \addplot[red, dashed] coordinates {(90,34.5) (0,34.5)};
    \addplot[red, dashed] coordinates {(90,0)    (90,34.5)};
    \node[red] at (axis cs: 75, 35.5) {\small $R \approx 34.5'$ bei $z=90^{\circ}$};
  \end{axis}
\end{tikzpicture}
\caption{Atmosphärische Refraktion $R(z)$ als Funktion des Zenitwinkels $z$
         nach der Formel von \citet{Meeus1991}.
         Am Horizont ($z = 90^{\circ}$) erreicht die Brechung ihr Maximum
         von $\approx \SI{34.5}{\arcminute}$,
         was dem scheinbaren Durchmesser der Sonne entspricht.}
\label{fig:refraktion_kurve}
\end{figure}

% ============================================================
\section{Dunkelheit des Universum}
\label{sec:olbers}
% ============================================================

\subsection{Fragestellung}

Das Sonnenlicht wird in der Erdatmosphäre harmonisch kaskadiert
gebrochen und erreicht die Erdoberfläche in biologisch optimaler
Form (Abschnitte~\ref{sec:kaskade}–\ref{sec:faecher}).
Doch trotz der Sonne und der Milliarden anderer Sterne
im Universum herrscht nachts Dunkelheit.
Dieses scheinbare Paradoxon – bekannt als \emph{Olbers-Paradoxon}
\citep{Harrison1987} – lässt sich vollständig formal auflösen.

\subsection{Das Paradoxon von Olbers (1823)}

Wäre das Universum unendlich groß, unendlich alt und
(gleichmäßig mit Sternen) gefüllt, so würde der Lichtfluss $F$
an der Erde von allen Richtungen gleich groß sein:
Jede Sichtlinie würde auf eine Sternoberfläche treffen,
und der Nachthimmel würde mit der Helligkeit einer Sonnenscheibe
leuchten.
Die Flächenhelligkeit $B$ wäre
\begin{equation}
  B = \frac{L_\odot}{4\pi R_\odot^2} \approx \SI{6.3e7}{\watt\per\square\metre\per\steradian},
  \label{eq:B_olbers}
\end{equation}
d.\ h., der gesamte Himmel würde leuchten wie die Sonnenscheibe.

\subsection{Formaler Nachweis: Lichtfluss einer Sternpopulation}

Betrachte eine homogene Sternpopulation mit Zahldichte
$n_\star$ [Sterne/m$^3$] und mittlerer Leuchtkraft $L_\star$.
Der differentielle Beitrag einer Kugelschale der Dicke $\dd r$
in Abstand $r$ zum Gesamtfluss ist:
\begin{equation}
  \dd F = n_\star\,L_\star\,\frac{1}{4\pi r^2}\cdot 4\pi r^2\,\dd r
        = n_\star\,L_\star\,\dd r.
  \label{eq:dF_schale}
\end{equation}
Das inverse quadratische Abstandsgesetz ($1/r^2$) wird durch das
quadratisch wachsende Schalenvolumen ($4\pi r^2\,\dd r$) genau
kompensiert.
Integration über alle Abstände ergibt im klassischen Modell
\begin{equation}
  F = \int_0^\infty n_\star\,L_\star\,\dd r \;\longrightarrow\; \infty.
  \label{eq:F_divergent}
\end{equation}
Dies ist der Kern des Paradoxons.

\subsection{Auflösung I: Endliches Alter und endliche Lichtlaufzeit}

Das Universum ist nicht unendlich alt, sondern entstand vor
$t_0 \approx \SI{13.8}{\giga\year}$ ($\approx \num{13.8}\,\mathrm{Gyr}$) im Urknall
\citep{Planck2018}.
Licht, das weiter als die Hubble-Radius $c\,t_0$ entfernt
entstand, hat die Erde noch nicht erreicht.
Das effektive Integrationsgebiet ist daher endlich:
\begin{equation}
  r_{\max} = c\,t_0
    \approx \SI{3e8}{\metre\per\second} \times \SI{4.35e17}{\second}
    \approx \SI{1.3e26}{\metre} \approx 13.8\,\mathrm{Gpc}.
  \label{eq:r_max}
\end{equation}
Der tatsächliche Nachtfluss beträgt dann
\begin{equation}
  F_{\mathrm{Nacht}} = n_\star\,L_\star\,r_{\max}
    \approx \SI{3e-9}{\watt\per\square\metre},
  \label{eq:F_nacht}
\end{equation}
für typische Werte $n_\star \approx 10^{8}\,\mathrm{Mpc}^{-3}$
und $L_\star \approx L_\odot$.
Dies entspricht der tatsächlich gemessenen Hintergrundstrahlung
des optischen Nachthimmels und ist um einen Faktor
$\sim 10^{22}$ kleiner als die Strahlungsleistung der Sonne
an der Erde ($\approx \SI{1361}{\watt\per\square\metre}$).

\subsection{Auflösung II: Kosmologische Rotverschiebung}

Das Universum expandiert; Licht weit entfernter Quellen wird
durch die kosmologische Rotverschiebung $z$ energieärmer:
\begin{equation}
  E_\gamma^{\mathrm{empfangen}} = \frac{E_\gamma^{\mathrm{emittiert}}}{1+z}.
  \label{eq:rotverschiebung}
\end{equation}
Für Quellen jenseits $z \gg 1$ wird selbst Röntgenstrahlung
ins Radioband verschoben und trägt nicht mehr zur optischen
Beleuchtung bei.
Der kosmologisch korrigierte Fluss fällt mit zunehmendem $z$
rasch ab:
\begin{equation}
  F(z) \propto \frac{L_\star}{(1+z)^4\,d_L^2(z)},
  \label{eq:F_korr}
\end{equation}
wobei $d_L(z)$ die Leuchtkraftdistanz ist, die für $z \to \infty$
faster als linear divergiert.
Die Kombination aus endlichem Weltalter
(Gleichung~\ref{eq:r_max}) und
kosmologischer Energiereduktion (Gleichung~\ref{eq:rotverschiebung})
bewirkt, dass das integrierte Licht aller Sterne und Galaxien
im Universum am Erdort auf einen winzigen Bruchteil der
Sonnenstrahlung reduziert wird.

\subsection{Auflösung III: Endliche Lebensdauer der Sterne}

Sterne brennen nicht ewig. Die Lebensdauer eines sonnenähnlichen
Sterns beträgt
\begin{equation}
  \tau_\star \approx \frac{M_\star\,\epsilon\,c^2}{L_\star}
             \approx \SI{10}{\giga\year} \approx 10\,\mathrm{Gyr},
  \label{eq:stern_lebensdauer}
\end{equation}
wobei $\epsilon \approx 0.007$ der Wirkungsgrad der
Wasserstofffusion ist.
Da $\tau_\star$ endlich ist, sind viele weiter entfernte
Sterne bereits erloschen, bevor ihr Licht die Erde erreicht.

\subsection{Warum das Sonnenlicht lokal perfekt, kosmisch jedoch
           unzureichend ist: Formaler Vergleich}

\begin{theorem}[Lokale Harmonie vs.\ kosmische Dunkelheit]
  Sei $\Phi_\odot$ der Strahlungsfluss der Sonne an der Erde,
  $\Phi_{\mathrm{Kosmos}}$ der integrierte Fluss aller
  übrigen Sterne und Galaxien im beobachtbaren Universum.
  Dann gilt
  \begin{equation}
    \frac{\Phi_{\mathrm{Kosmos}}}{\Phi_\odot}
    \approx \frac{n_\star\,L_\star\,r_{\max}}{L_\odot / (4\pi\,d_{\mathrm{SE}}^2)}
    \approx \num{2e-9}.
    \label{eq:verhaeltnis}
  \end{equation}
  Das Licht aller anderen Sterne zusammen ist also um neun
  Größenordnungen schwächer als das Sonnenlicht an der Erde.
\end{theorem}

\begin{proof}
  Einsetzen von $n_\star = 1.4 \times 10^{8}\,\mathrm{Mpc}^{-3}$
  \citep{Conselice2016},
  $L_\star = L_\odot = \SI{3.83e26}{\watt}$,
  $r_{\max} = c\,t_0 \approx \SI{4.26e26}{\metre}$ und
  $d_{\mathrm{SE}} = \SI{1.496e11}{\metre}$:
  \begin{align*}
    \Phi_{\mathrm{Kosmos}} &= n_\star\,L_\star\,r_{\max}
      = 1.4\times 10^{8}\,\mathrm{Mpc}^{-3}\times L_\odot
        \times \SI{4.26e26}{\metre}\\
      &\approx \SI{2.9e-9}{\watt\per\square\metre},\\
    \Phi_\odot &= \frac{L_\odot}{4\pi\,d_{\mathrm{SE}}^2}
      = \frac{\SI{3.83e26}{\watt}}{4\pi (\SI{1.496e11}{\metre})^2}
      \approx \SI{1361}{\watt\per\square\metre}.
  \end{align*}
  Damit ergibt sich
  $\Phi_{\mathrm{Kosmos}}/\Phi_\odot \approx 2.1\times 10^{-9}$,
  was die Ungleichung~\eqref{eq:verhaeltnis} bestätigt.
\end{proof}

\subsection{Zusammenfassung: Warum das Universum dunkel ist}

Drei voneinander unabhängige, aber gleichzeitig wirkende
Ursachen erklären die kosmische Dunkelheit:

\begin{enumerate}
  \item \textbf{Endliches Weltalter:}
        Das Licht weit entfernter Sterne hat die Erde noch nicht
        erreicht (Gleichung~\ref{eq:r_max}).
        Der sichtbare Kosmos ist endlich.

  \item \textbf{Kosmologische Rotverschiebung:}
        Die Expansion des Universums streckt die Wellenlänge
        des Lichts und reduziert seine Energie
        (Gleichung~\ref{eq:rotverschiebung}).
        Weit entferntes Licht ist für das menschliche Auge
        unsichtbar geworden.

  \item \textbf{Endliche Sternlebensdauer:}
        Viele Sterne jenseits der kosmologischen Reichweite
        sind bereits erloschen (Gleichung~\ref{eq:stern_lebensdauer}).
\end{enumerate}

Das Sonnenlicht hingegen wirkt durch die kaskadierende,
harmonische Schichtbrechung der Erdatmosphäre exakt so,
dass es auf der Erdoberfläche in biologisch optimaler Intensität,
Spektralverteilung und Richtungscharakteristik ankommt –
ein Zusammenspiel von Schichtbrechung und astronomischer Entfernung,
das im gesamten bekannten Universum so nur an der Erde nachgewiesen ist
(vgl. Gleichung~\ref{eq:verhaeltnis}).

% ============================================================
\section{Schluss}
\label{sec:schluss}
% ============================================================

Die vorliegende Analyse zeigt, dass die schichtweise aufgebaute
Erdatmosphäre die Ausbreitung des Sonnenlichts beim Auf- und
Untergang in mehrfacher Hinsicht beeinflusst:

\begin{enumerate}
  \item \textbf{Kontinuierliche Krümmung der Strahlbahn.}
        Jede Atmosphärenschicht mit $\dd n/\dd h < 0$ biegt den
        Lichtstrahl um einen kleinen Betrag zur Erde hin.
        Die Summe aller Beiträge ergibt am Horizont
        $R \approx \SI{34.5}{\arcminute}$.

  \item \textbf{Verlängerte Tagesdauer.}
        Der Lichtstrahl erreicht den Beobachter,
        obwohl die Sonne geometrisch bereits ($\approx \SI{34.5}{\arcminute}$
        $\approx$ einem Sonnendurchmesser) unter dem Horizont steht.
        Dies entspricht einer Zeitverschiebung von ca.\ \SI{2}{\minute}
        \SI{20}{\second} pro Horizont.

  \item \textbf{Strahlenförmige Lichtverteilung vom Sonnenzentrum.}
        Da alle Lichtstrahlen von der Sonnenscheibe ausgehen
        und nahezu parallel sind, erzeugt die perspektivische
        Projektion den Eindruck, sie strömten strahlenförmig
        vom scheinbaren Sonnenmittelpunkt aus –
        ein rein geometrischer Projektionseffekt in Kombination
        mit der schichtweisen Brechung
        (Gleichungen~\ref{eq:delta_R}–\ref{eq:fluchtpunkt}).

  \item \textbf{Dominanz der Troposphäre.}
        Die untersten $\SI{12}{\kilo\metre}$ der Atmosphäre
        liefern $\approx 81\%$ der Gesamtbrechung,
        da hier der stärkste Brechungsindexgradient und die
        größte Luftdichte vorliegen.

  \item \textbf{Kosmische Dunkelheit trotz Milliarden von Sternen.}
        Das Olbers-Paradoxon wird durch das endliche Weltalter,
        die kosmologische Rotverschiebung und die endliche
        Sternlebensdauer aufgelöst.
        Das integrierte Licht aller anderen Sterne ist um
        $\approx 10^9$ schwächer als das Sonnenlicht an der Erde
        (Gleichung~\ref{eq:verhaeltnis}).
\end{enumerate}

Die berechneten Werte stimmen gut mit astronomischen Beobachtungen überein
\citep{Stone1996} und bilden die Grundlage für
Sonnenuhr-Korrekturen, nautische Almanache und die Modellierung
von Fernsehrundfunk-Signalausbreitung über den Horizont hinaus.

\subsection{Schlussfolgerungen}

Es wurde gezeigt, dass die Erdatmosphäre mit ihren konzentrischen
Schichten unterschiedlicher Dichte und Brechungsindex das Sonnenlicht
beim Auf- und Untergang gemäß dem Snellschen Brechungsgesetz
(in seiner sphärischen Form, dem Bouguer-Gesetz) stetig krümmt.
Die mathematische Herleitung über die Strahlgleichung in
sphärischer Geometrie liefert eine Gesamtablenkung von
$R_{\mathrm{Horizont}} \approx \SI{34.5}{\arcminute}$,
wobei die Troposphäre den weitaus größten Beitrag leistet.

Als Folge dieser schichtweisen Brechung:
\begin{itemize}
  \item ist die Sonne beim Auf- und Untergang für den irdischen
        Beobachter sichtbar, obwohl sie geometrisch un\-ter dem
        Horizont liegt;
  \item erscheint das Licht der Sonnenscheibe strahlenförmig
        vom scheinbaren Sonnenzentrum aus über den Horizont
        verteilt; dies ist der Satz der krepuskulären Strahlen,
        der durch die perspektivische Konvergenz nahezu paralleler
        Lichtstrahlen auf den Fluchtpunkt Sonne entsteht.
\end{itemize}

Beide Effekte sind vollständig durch die klassische geometrische Optik
und das Brechungsindexprofil der Standardatmosphäre erklärbar
und stellen keinen empirisch neuen Befund dar.
Sie sind in ihrer Gesamtheit das unmittelbare Ergebnis des
Grundprinzips: \emph{Licht folgt dem Weg, der die optische
Weglänge extremiert – und die Atmosphäre sorgt dafür, dass dieser
Weg zur Erde hin gebogen ist.}

Darüber hinaus wurde formal nachgewiesen (Abschnitt~\ref{sec:olbers}),
dass das Sonnenlicht nicht das gesamte Universum erleuchten kann.
Drei physikalische Ursachen – das endliche Weltalter, die
kosmologische Rotverschiebung und die endliche Sternlebensdauer –
sorgen dafür, dass das integrierte kosmische Hintergrundlicht
um neun Größenordnungen schwächer ist als die Sonneneinstrahlung
an der Erde.
Die harmonische kaskadierende Schichtbrechung der Erdatmosphäre
ist damit ein lokal streng optimierter Prozess, dessen Wirkung
auf den winzigen Bereich zwischen Atmosphäre und Erdoberfläche
beschränkt bleibt und keinerlei kosmologische Reichweite besitzt.

\subsection{Ausblick}

Die vorliegende Arbeit legt ein konsistentes physikalisches Fundament
für das Verständnis der atmosphärischen Lichtbrechung und ihrer
kosmologischen Einbettung.
Gleichwohl eröffnen sich zahlreiche weiterführende Fragestellungen,
die sowohl grundlagenphysikalischer als auch angewandter Natur sind.
Die folgenden Abschnitte skizzieren die bedeutendsten offenen
Forschungsrichtungen und ihre Relevanz für die jeweiligen wissenschaftlichen
Gemeinschaften.

\subsubsection{Präzisionsmodellierung des Brechungsindexprofils}

Das in dieser Arbeit verwendete Exponentialprofil $n(h) \approx 1 +
\delta_0\,e^{-h/H}$ (Gleichung~\ref{eq:n_exp}) ist eine gut motivierte
Näherung für die Standardatmosphäre, erfasst jedoch nicht die
ausgeprägte Temperaturinversion in der Stratosphäre (Ozonschicht,
$\approx\SI{25}{\kilo\metre}$) noch die turbulenten Grenzschicht\-prozesse
in der untersten Troposphäre.
Ein lohnenswerter Folgeschritt wäre die Implementierung des
NRLMSISE-00-Atmosphärenmodells \citep{Picone2002} als tabelliertes
Referenzprofil, das saisonale, breitengrad- und tageszeit\-abhängige
Variationen von $p(h)$, $T(h)$ und der Zusammensetzung beinhaltet.
Insbesondere die Dispersion -- die Wellenlängenabhängigkeit von
$n(\lambda, h)$ -- sollte für eine vollständige chromatische
Analyse einbezogen werden:
\begin{equation}
  n(\lambda,h) - 1
  = \frac{K(\lambda)\,p(h)}{T(h)},
  \qquad
  K(\lambda) = K_0\!\left(1 + \frac{b}{\lambda^2}\right),
  \label{eq:dispersion}
\end{equation}
wobei $b \approx \SI{1.36e-14}{\metre\squared}$ der
Cauchy-Koeffizient ist \citep{BornWolf1999}.
Dies würde eine direkte Berechnung der atmosphärischen Dispersion
-- der unterschiedlichen Brechung von Rot- und Blaulicht -- erlauben
und das Phänomen des grünen Blitzes (\emph{green flash}) quantitativ
erklären, das nur unter Präzisionsbedingungen beobachtbar ist
\citep{Greenler1980}.

\subsubsection{Numerische Raytracing-Simulation in drei Dimensionen}

Die in dieser Arbeit entwickelte analytische Strahlgleichung
(Gleichung~\ref{eq:R_integral}) setzt sphärische Symmetrie und
einen glatten Brechungsindexgradienten voraus.
Reale atmosphärische Bedingungen weichen von dieser Idealisierung ab:
horizontale Brechungsindexgradienten durch Temperaturinversionen,
Luftpakete unterschiedlicher Feuchte und turbulente Strömungen
erzeugen Fluktuationen im Strahlverlauf.
Ein dreidimensionales Raytracing-Modell auf Basis der Hamiltonschen
Strahlgleichungen
\begin{equation}
  \dot{\mathbf{r}} = \frac{\partial \mathcal{H}}{\partial \mathbf{k}}, \qquad
  \dot{\mathbf{k}} = -\frac{\partial \mathcal{H}}{\partial \mathbf{r}},
  \qquad \mathcal{H} = n(\mathbf{r}) - |\mathbf{k}|,
  \label{eq:hamilton_ray}
\end{equation}
– numerisch gelöst über eine adaptiven Runge-Kutta-Integration –
würde die Strahlbahn unter realistischen Bedingungen simulieren
und könnte mit Lidar- und Radiooccultations-Messungen
(z.\,B. GPS-RO-Profilen von COSMIC-2 \citep{Anthes2008}) validiert werden.
Eine solche Simulation hätte direkte Anwendungen in der
Präzisionsastrometrie (Korrekturen für optische Teleskope nahe
dem Horizont), der Satellitennavigation (GPS-Troposphärenkorrektur)
und der atmosphärischen Fernerkundung.

\subsubsection{Harmonische Analyse der kaskadierenden Brechungsfolge}

Die vorliegende Arbeit charakterisiert die Folge der Brechungswinkel
$\delta_1, \delta_2, \ldots, \delta_N$ als aufsummierenden Prozess
(Gleichung~\ref{eq:R_kaskade}).
Von wissenschaftlichem Interesse ist die Frage, ob diese Folge
eine Harmonizitätseigenschaft im funktionalanalytischen Sinne besitzt.
Konkret ist zu untersuchen, ob die Funktion
\begin{equation}
  f(h) \;=\; -\frac{1}{n(h)}\,\dv{n}{h}
           \;=\; \frac{\delta_0/H\,e^{-h/H}}{1 + \delta_0\,e^{-h/H}}
  \label{eq:harmonik_f}
\end{equation}
im Sinne von
\begin{equation}
  \int_0^\infty f(h)\,\dd h = \delta_0
  \label{eq:harmonik_integral}
\end{equation}
eine optimale Verteilung der Brechungsbeiträge auf die Atmosphärenschichten
darstellt, die -- verglichen mit anderen denkbaren Profilen $n'(h)$ mit
gleichem $n(0)$ und $n(\infty)=1$ -- den mittleren quadratischen
Abbildungsfehler für polychromatisches Licht minimiert.
Diese Optimierungsperspektive würde das in der vorliegenden Arbeit
verwendete Konzept der ``harmonischen Brechung'' mathematisch auf ein
festes Fundament stellen und könnte Anknüpfungspunkte zur
Variationsrechnung (Fermat-Prinzip in inhomogenen Medien) liefern.

\subsubsection{Klimawandelbedingte Veränderungen der atmosphärischen Refraktion}

Der anthropogene Klimawandel verändert das Temperatur- und
Feuchtigkeitsprofil der Atmosphäre:
Die Troposphäre erwärmt sich ($+\SI{0.2}{\kelvin}$ pro Dekade),
die Tropopause senkt sich (durch stratosphärische Abkühlung),
und der Wasserdampfgehalt steigt entsprechend dem
Clausius-Clapeyron-Gesetz um $\approx 7\%$ pro Kelvin
Erwärmung.
Da der Brechungsindex der feuchten Luft vom Wasserdampfpartialdruck $e$
abhängt,
\begin{equation}
  n_{\mathrm{feucht}}(h) - 1
  \;=\; \frac{77.6}{T}\,p + \frac{3.73\times10^5}{T^2}\,e,
  \label{eq:n_feucht}
\end{equation}
(Smith-Weintraub-Formel \citep{SmithWeintraub1953}),
verändert sich das effektive Skalenhöhenprofil $H_{\mathrm{eff}}(t)$
mit der Zeit.
Für die Astronomie und die Geodäsie bedeutet dies, dass die
klassischen Refraktionstafeln (basierend auf der
ICAO-Standardatmosphäre von 1964) im späten 21.\,Jahrhundert
regelmäßig neu kalibriert werden müssten.
Eine quantitative Abschätzung dieser Drift wäre ein konkreter
Beitrag zur Klimafolgenforschung in der Metrologie.

\subsubsection{Biologische Relevanz der kaskadierenden Schichtbrechung}

Die kaskadierende Schichtbrechung bewirkt nicht nur eine geometrische
Ablenkung des Lichtstrahls, sondern auch eine spektrale Modifikation
der solaren Bestrahlungsstärke $E(\lambda)$ an der Erdoberfläche.
Durch Rayleigh-Streuung ($ \propto \lambda^{-4}$) und die Ozonabsorption
im UV-Bereich (Hartley-Band, $200\text{--}\SI{300}{\nano\metre}$)
hat die Atmosphäre ein biologisch optimales Transmissionsfenster
im Bereich $\SI{400}{\nano\metre}$ bis $\SI{700}{\nano\metre}$
(photosynthetisch aktive Strahlung, PAR) geformt.
Eine interessante offene Frage ist, ob die mehrfache Brechung an
Schichtgrenzen diese spektrale Selektivität zusätzlich beeinflusst --
etwa durch winkelabhängige Fresnel-Reflexionsverluste
\begin{equation}
  R_{\mathrm{Fresnel}}(\theta) =
  \left(\frac{n_i \cos\theta_i - n_{i+1}\cos\theta_{i+1}}
             {n_i \cos\theta_i + n_{i+1}\cos\theta_{i+1}}\right)^2,
  \label{eq:fresnel}
\end{equation}
die für flache Einfallswinkel signifikant werden könnten.
Die quantitative Abschätzung dieser Verluste und ihrer spektralen
Verteilung wäre für die Modellierung des solaren Energiehaushalts
der Erde und für Photovoltaik-Auslegungen in Horizontnähe von
praktischem Interesse.

\subsubsection{Vergleich mit anderen Planeten des Sonnensystems}

Ein natürliches Labor für die Theorie der kaskadierenden Schichtbrechung
sind die Atmosphären anderer Planeten.
Die Venus, Mars, Jupiter und Saturn besitzen Atmosphären
mit völlig anderen Dichte-, Temperatur- und Zusammensetzungsprofilen.
Für die Mars-Atmosphäre (überwiegend CO$_2$, $p_0 \approx \SI{6}{\hecto\pascal}$,
$\delta_0 \approx 5\times10^{-5}$) beträgt die Horizontalrefraktion
nur $\approx \SI{6}{\arcminute}$ \citep{Kliore1972},
was einer biologisch-atmosphärischen Schutzfunktion wie auf der Erde
nicht nahekommt.
Im Gegensatz dazu besitzt die Venus eine so dichte Atmosphäre
($p_0 \approx \SI{93}{bar}$), dass totale Ãœberrefraktion auftreten kann:
Licht kann um die Venusoberfläche herumgeführt werden,
was von Sondenstartfenstern und Landeplanung berücksichtigt werden muss.
Ein systematischer Vergleich
\begin{equation}
  R_{\mathrm{Horizont}}^{(\mathrm{Planet})}
  = f\!\left(\delta_0^{(\mathrm{Planet})},\, H^{(\mathrm{Planet})},\,
             r_E^{(\mathrm{Planet})}\right)
  \label{eq:R_planeten}
\end{equation}
über das Gleichungssystem~\eqref{eq:R_integral} würde nicht nur das
Verständnis exoplanetarer Atmosphären vertiefen,
sondern auch Hinweise auf die Bedingungen liefern,
unter denen Leben -- angewiesen auf eine mäßig abgeschwächte,
harmonisch gebrochene Sonnenstrahlung -- entstehen und überleben kann.

\subsubsection{Verbindung zur Kosmologie: Atmosphärische Refraktion
               als analoges Modell für Gravitationslinsen}

Die formale Analogie zwischen der atmosphärischen Strahlbrechung
in einem sphärisch symmetrischen Medium (Bouguer-Gesetz,
Gleichung~\ref{eq:bouguer}) und der Ablenkung von Lichtstrahlen
in einem sphärisch symmetrischen Gravitationsfeld
(Schwarzschild-Metrik) ist wohlbekannt \citep{SchneiderEhlers1992}:
\begin{equation}
  \alpha_{\mathrm{grav}} = \frac{4GM}{c^2 b},
  \qquad \text{(Einstein-Ablenkung)}
  \label{eq:einstein}
\end{equation}
wobei $b$ der Stoßparameter und $M$ die ablenkende Masse ist.
Die vorliegende Arbeit legt nahe, dass das atmosphärische
Brechungsmodell -- insbesondere die numerische Integration
von Gleichung~\eqref{eq:R_integral} -- als Testbed für
Gravitationslinsen-Algorithmen dienen könnte:
Durch Wahl von $n(r) = 1 + 2GM/(c^2 r)$ reproduziert das
Bouguer-Integral exakt die Einsteinsche Ablenkung im Schwachfeldlimit.
Damit liefert die Atmosphärenoptik ein anschauliches und labortaugliches
Analogon zur allgemeinen Relativitätstheorie, das in
Lehre und Forschung gleichermaßen fruchtbar eingesetzt werden kann.

\subsubsection{Weiterentwicklung der kaskadierenden Brechung
               zu einem vollständigen Transfermodell}

Die in dieser Arbeit eingeführte Darstellung der Atmosphäre als
endliche Schichtfolge (Gleichung~\ref{eq:R_kaskade}) lässt sich
zu einem vollständigen Strahlungs\-transfermodell erweitern,
das neben Refraktion auch Absorption, Streuung und Emission
der Atmosphärenschichten berücksichtigt.
Das Beer-Lambert-Gesetz für die Transmission durch eine Schicht $i$
\begin{equation}
  \tau_i(\lambda) = \exp\!\left(-\int_{h_i}^{h_{i+1}}
    \sigma_{\mathrm{ext}}(\lambda,h)\,\rho(h)\,\dd h \right)
  \label{eq:beer_lambert}
\end{equation}
kombiniert mit der geometrischen Pfadverlängerung durch die Brechung
(Airmass-Faktor $X(z) \approx \sec z + \text{Korrekturterme}$)
würde die Grundlage für hochpräzise photometrische Korrekturen
in der bodengebundenen Astronomie und für die Modellierung des
Klimastrahlungsantriebs (Forcing) liefern.
Insbesondere die Kopplung mit dreidimensionalen Klimamodellen
(z.\,B. ECMWF-IFS oder NCAR-CESM) wäre ein bedeutender Schritt
hin zu einem integrierten Erd\-system\-modell,
das Lichtausbreitung, Schichtbrechung und biologische
Wirkung des Sonnenlichts in einem konsistenten Rahmen behandelt.

\subsubsection{Abschließende Perspektive}

Die kaskadierende, harmonische Brechung des Sonnenlichts in den
Schichten der Erdatmosphäre erweist sich bei näherer Betrachtung
als ein Phänomen von außerordentlicher Tiefe:
Es verbindet die klassische geometrische Optik (Snell, Bouguer, Fermat)
mit der modernen Atmosphärenphysik, der Kosmologie (Olbers-Paradoxon,
Gravitationslinsen) und den Lebenswissenschaften (biologisches
Strahlungsfenster, Tageslichtrhythmus der Biosphäre).
Das Zusammenspiel aus dem enormen Größenverhältnis Sonne/Erde,
der endlichen Entfernung und dem präzise gestuften Brechungsindexprofil
der Atmosphäre erzeugt eine lokale Lichtqualität,
die weder auf dem Mond (keine Atmosphäre, keine Brechung)
noch auf der Venus (Überrefraktion, hochgiftige Atmosphäre)
noch auf dem Mars (zu dünne Atmosphäre, zu geringe Brechung)
in vergleichbarer Weise auftritt.
In diesem Sinne ist die harmonische Schichtbrechung der Erdatmosphäre
nicht nur ein ästhetisches Phänomen des Sonnenauf- und -untergangs,
sondern eine notwendige Vorbedingung für die Existenz komplexen
Lebens auf der Erdoberfläche --
eine Aussage, die durch die in dieser Arbeit vorgelegten formalen
Nachweise eine feste quantitative Grundlage erhält.

% ============================================================
%  Literaturverzeichnis
% ============================================================
\bibliographystyle{plainnat}
\begin{thebibliography}{99}

\bibitem[Born \& Wolf(1999)]{BornWolf1999}
Born, M., \& Wolf, E. (1999).
\newblock \textit{Principles of Optics} (7th ed.).
\newblock Cambridge University Press.
\newblock \doi{10.1017/CBO9781139644181}

\bibitem[Greenler(1980)]{Greenler1980}
Greenler, R. (1980).
\newblock \textit{Rainbows, Halos, and Glories}.
\newblock Cambridge University Press.

\bibitem[Hohenkerk \& Sinclair(1985)]{Hohenkerk1988}
Hohenkerk, C.\,Y., \& Sinclair, A.\,T. (1985).
\newblock The computation of angular atmospheric refraction at large
  zenith angles.
\newblock \textit{H.M. Nautical Almanac Office Technical Note}, 63.

\bibitem[Meeus(1991)]{Meeus1991}
Meeus, J. (1991).
\newblock \textit{Astronomical Algorithms}.
\newblock Willmann-Bell, Richmond, Virginia.

\bibitem[Minnaert(1954)]{MinnaertNaturalOptics}
Minnaert, M. (1954).
\newblock \textit{The Nature of Light and Colour in the Open Air}.
\newblock Dover Publications, New York.

\bibitem[Penndorf(1957)]{PenndorfRefractivity1957}
Penndorf, R. (1957).
\newblock Tables of refractive index for standard air and the Rayleigh
  scattering coefficient for the spectral region between 0.2 and 20.0 $\mu$.
\newblock \textit{Journal of the Optical Society of America}, 47(2), 176--182.
\newblock \doi{10.1364/JOSA.47.000176}

\bibitem[Ramsauer(1932)]{Ramsauer1932}
Ramsauer, C. (1932).
\newblock Zur Theorie der atmosphärischen Refraktion.
\newblock \textit{Annalen der Physik}, 404(4), 425--448.

\bibitem[Smart(1977)]{Smart1977}
Smart, W.\,M. (1977).
\newblock \textit{Textbook on Spherical Astronomy} (6th ed., revised by Green).
\newblock Cambridge University Press.

\bibitem[Stone(1996)]{Stone1996}
Stone, R.\,C. (1996).
\newblock An accurate method for computing atmospheric refraction.
\newblock \textit{Publications of the Astronomical Society of the Pacific},
  108, 1051--1058.
\newblock \doi{10.1086/133831}

\bibitem[Conselice et al.(2016)]{Conselice2016}
Conselice, C.\,J., Wilkinson, A., Duncan, K., \& Mortlock, A. (2016).
\newblock The evolution of galaxy number density at $z < 8$ and its
  implications.
\newblock \textit{The Astrophysical Journal}, 830(2), 83.
\newblock \doi{10.3847/0004-637X/830/2/83}

\bibitem[Harrison(1987)]{Harrison1987}
Harrison, E.\,R. (1987).
\newblock \textit{Darkness at Night: A Riddle of the Universe}.
\newblock Harvard University Press, Cambridge, MA.

\bibitem[Planck Collaboration(2018)]{Planck2018}
Planck Collaboration (2018).
\newblock Planck 2018 results. VI. Cosmological parameters.
\newblock \textit{Astronomy \& Astrophysics}, 641, A6.
\newblock \doi{10.1051/0004-6361/201833910}

\bibitem[Anthes et al.(2008)]{Anthes2008}
Anthes, R.\,A., et al. (2008).
\newblock The COSMIC/FORMOSAT-3 mission: Early results.
\newblock \textit{Bulletin of the American Meteorological Society}, 89(3),
  313--333.
\newblock \doi{10.1175/BAMS-89-3-313}

\bibitem[Kliore et al.(1972)]{Kliore1972}
Kliore, A., Fjeldbo, G., Seidel, B.\,L., \& Rasool, S.\,I. (1972).
\newblock Mariners 6 and 7: Radio occultation measurements of the atmosphere
  of Mars.
\newblock \textit{Journal of the Atmospheric Sciences}, 29(3), 457--474.

\bibitem[Picone et al.(2002)]{Picone2002}
Picone, J.\,M., Hedin, A.\,E., Drob, D.\,P., \& Aikin, A.\,C. (2002).
\newblock NRLMSISE-00 empirical model of the atmosphere.
\newblock \textit{Journal of Geophysical Research: Space Physics}, 107(A12),
  1468.
\newblock \doi{10.1029/2002JA009430}

\bibitem[Schneider \& Ehlers(1992)]{SchneiderEhlers1992}
Schneider, P., Ehlers, J., \& Falco, E.\,E. (1992).
\newblock \textit{Gravitational Lenses}.
\newblock Springer-Verlag, Berlin.
\newblock \doi{10.1007/978-3-662-03758-4}

\bibitem[Smith \& Weintraub(1953)]{SmithWeintraub1953}
Smith, E.\,K., \& Weintraub, S. (1953).
\newblock The constants in the equation for atmospheric refractive index at
  radio frequencies.
\newblock \textit{Proceedings of the IRE}, 41(8), 1035--1037.
\newblock \doi{10.1109/JRPROC.1953.274297}

% ── 17 ergänzte Quellen ────────────────────────────────────────────────────

\bibitem[Allen(1976)]{Allen1976}
Allen, C.\,W. (1976).
\newblock \textit{Astrophysical Quantities} (3rd ed.).
\newblock Athlone Press, London.

\bibitem[Auer \& Standish(2000)]{AuerStandish2000}
Auer, L.\,H., \& Standish, E.\,M. (2000).
\newblock Astronomical refraction: Computational method for all zenith angles.
\newblock \textit{The Astronomical Journal}, 119(5), 2472--2474.
\newblock \doi{10.1086/301325}

\bibitem[Baldini et al.(2020)]{Baldini2020}
Baldini, L., Bindi, V., Corti, V., De Mitri, I., et al. (2020).
\newblock Spectral characterization of the photosynthetically active radiation
  and its atmospheric transmission at sea level.
\newblock \textit{Journal of Atmospheric and Solar-Terrestrial Physics}, 207,
  105364.
\newblock \doi{10.1016/j.jastp.2020.105364}

\bibitem[Bean \& Dutton(1966)]{BeanDutton1966}
Bean, B.\,R., \& Dutton, E.\,J. (1966).
\newblock \textit{Radio Meteorology}.
\newblock National Bureau of Standards Monograph 92,
  U.S. Government Printing Office, Washington D.C.

\bibitem[Chandrasekar(1960)]{Chandrasekar1960}
Chandrasekhar, S. (1960).
\newblock \textit{Radiative Transfer}.
\newblock Dover Publications, New York.

\bibitem[Edlén(1966)]{Edlen1966}
Edlén, B. (1966).
\newblock The refractive index of air.
\newblock \textit{Metrologia}, 2(2), 71--80.
\newblock \doi{10.1088/0026-1394/2/2/002}

\bibitem[Filippenko(1982)]{Filippenko1982}
Filippenko, A.\,V. (1982).
\newblock The importance of atmospheric differential refraction in
  spectrophotometry.
\newblock \textit{Publications of the Astronomical Society of the Pacific},
  94(554), 715--721.
\newblock \doi{10.1086/131052}

\bibitem[Garfinkel(1967)]{Garfinkel1967}
Garfinkel, B. (1967).
\newblock Astronomical refraction in a polytropic atmosphere.
\newblock \textit{The Astronomical Journal}, 72, 235--254.
\newblock \doi{10.1086/110225}

\bibitem[Hecht(2017)]{Hecht2017}
Hecht, E. (2017).
\newblock \textit{Optics} (5th ed.).
\newblock Pearson Education, Harlow.

\bibitem[ICAO(1993)]{ICAO1993}
International Civil Aviation Organization (1993).
\newblock \textit{Manual of the ICAO Standard Atmosphere} (3rd ed.).
\newblock ICAO Doc 7488/3, Montreal.

\bibitem[Kovalevsky \& Seidelmann(2004)]{Kovalevsky2004}
Kovalevsky, J., \& Seidelmann, P.\,K. (2004).
\newblock \textit{Fundamentals of Astrometry}.
\newblock Cambridge University Press.
\newblock \doi{10.1017/CBO9781139106832}

\bibitem[Liou(2002)]{Liou2002}
Liou, K.-N. (2002).
\newblock \textit{An Introduction to Atmospheric Radiation} (2nd ed.).
\newblock Academic Press, San Diego.

\bibitem[Owens(1967)]{Owens1967}
Owens, J.\,C. (1967).
\newblock Optical refractive index of air: Dependence on pressure, temperature
  and composition.
\newblock \textit{Applied Optics}, 6(1), 51--59.
\newblock \doi{10.1364/AO.6.000051}

\bibitem[Peck \& Reeder(1972)]{PeckReeder1972}
Peck, E.\,R., \& Reeder, K. (1972).
\newblock Dispersion of air.
\newblock \textit{Journal of the Optical Society of America}, 62(8), 958--962.
\newblock \doi{10.1364/JOSA.62.000958}

\bibitem[Rees(1989)]{Rees1989}
Rees, M.\,J. (1989).
\newblock \textit{The New Physics} (Chapter: The Expanding Universe).
\newblock Cambridge University Press.

\bibitem[Seidelmann(1992)]{Seidelmann1992}
Seidelmann, P.\,K. (Ed.) (1992).
\newblock \textit{Explanatory Supplement to the Astronomical Almanac}.
\newblock University Science Books, Mill Valley, California.

\bibitem[Thayer(1974)]{Thayer1974}
Thayer, G.\,D. (1974).
\newblock An improved equation for the radio refractive index of air.
\newblock \textit{Radio Science}, 9(10), 803--807.
\newblock \doi{10.1029/RS009i010p00803}

\end{thebibliography}

\end{document}

